%! AP Statistics — Unit 4, Lesson 4.8 — Group Activity (Answer Key)
\documentclass[10pt]{article}
\usepackage{apstats-article}
\usepackage{apstats-key}
\newcommand{\TallMath}[1]{$\displaystyle #1\rule[-1.4em]{0pt}{3.2em}$}

\begin{document}

\pageheader{Unit 4, Lesson 4.8}{Group Activity: Expected Value and Spread --- Answer Key}
\namepartnerperiod

\vspace{0.1in}

\begin{tcolorbox}[colback=sky!60, colframe=navy, title=\textbf{Tier R --- Remediate},
    arc=2mm, boxrule=0.8pt, left=4mm, right=4mm, top=3mm, bottom=3mm]
\small
$X$ = number of text messages a student sends per class period.
\begin{center}
\renewcommand{\arraystretch}{1.4}
\begin{tabular}{@{}ccccc@{}}
$x$ & 0 & 1 & 2 & 3 \\
\hline
$P(X=x)$ & 0.40 & 0.35 & 0.15 & 0.10 \\
\end{tabular}
\end{center}

\renewcommand{\arraystretch}{1.9}
\begin{tabularx}{\linewidth}{@{} c c X @{}}
$x$ & $P(X=x)$ & $x \cdot P(X=x)$ \\
\hline
0 & 0.40 & \ans{0.00} \\
1 & 0.35 & \ans{0.35} \\
2 & 0.15 & \ans{0.30} \\
3 & 0.10 & \ans{0.30} \\
\hline
\textbf{Sum} & & $\mu_X =$ \ans{0.95} \\
\end{tabularx}

\begin{enumerate}[leftmargin=1.5em, itemsep=8pt]
  \item Fill in the computation table and find $\mu_X$.

  \item Interpret $\mu_X$ in context.\\
        \ansline{Over many class periods, students send an average of about}\\[1pt]\ansline{0.95 text messages per class period.}

  \item Is $\mu_X$ a possible value of $X$? Explain.\\
        \ansline{No. $X$ takes whole numbers only; 0.95 is a long-run average.}
\end{enumerate}
\end{tcolorbox}

\vspace{0.12in}

\begin{tcolorbox}[colback=goldbg!60, colframe=goldacc, title=\textbf{Tier A --- Approaching Proficiency},
    arc=2mm, boxrule=0.8pt, left=4mm, right=4mm, top=3mm, bottom=3mm]
\small
$Y$ = number of hours a student exercises per week.
\begin{center}
\renewcommand{\arraystretch}{1.4}
\begin{tabular}{@{}ccccc@{}}
$y$ & 0 & 1 & 2 & 3 \\
\hline
$P(Y=y)$ & 0.20 & 0.35 & 0.30 & 0.15 \\
\end{tabular}
\end{center}

\renewcommand{\arraystretch}{2.0}
\begin{tabularx}{\linewidth}{@{} c c c c X @{}}
$y$ & $P$ & $y \cdot P$ & $(y - \mu_Y)^2$ & $(y-\mu_Y)^2 \cdot P$ \\
\hline
0 & 0.20 & \ans{0.00} & \ans{1.96} & \ans{0.392} \\
1 & 0.35 & \ans{0.35} & \ans{0.16} & \ans{0.056} \\
2 & 0.30 & \ans{0.60} & \ans{0.36} & \ans{0.108} \\
3 & 0.15 & \ans{0.45} & \ans{2.56} & \ans{0.384} \\
\hline
& & $\mu_Y=$ \ans{1.40} & & $\sigma^2_Y=$ \ans{0.940} \\
\end{tabularx}

\vspace{4pt}
$\sigma_Y = $ \ans{$\sqrt{0.940} \approx 0.970$ hours}

\begin{enumerate}[leftmargin=1.5em, itemsep=8pt]
  \item Interpret $\mu_Y$ in context.\\
        \ansline{Over many randomly selected students, the average amount of}\\[1pt]\ansline{exercise is about 1.40 hours per week.}

  \item Interpret $\sigma_Y$ in context.\\
        \ansline{Hours of exercise per week typically vary from the mean of}\\[1pt]\ansline{1.40 hours by about 0.97 hours.}
\end{enumerate}
\end{tcolorbox}

\vspace{0.12in}

\begin{tcolorbox}[colback=greenbg!60, colframe=greenacc, title=\textbf{Tier E --- Extension},
    arc=2mm, boxrule=0.8pt, left=4mm, right=4mm, top=3mm, bottom=3mm]
\small
Two carnival games each cost \$3 to play.

\textbf{Game A:}
\begin{center}
\renewcommand{\arraystretch}{1.3}
\begin{tabular}{@{}cccc@{}}
Prize ($x$) & \$0 & \$3 & \$10 \\
\hline
$P(X=x)$ & 0.60 & 0.30 & 0.10 \\
\end{tabular}
\end{center}

\textbf{Game B:}
\begin{center}
\renewcommand{\arraystretch}{1.3}
\begin{tabular}{@{}cccc@{}}
Prize ($y$) & \$0 & \$2 & \$6 \\
\hline
$P(Y=y)$ & 0.40 & 0.40 & 0.20 \\
\end{tabular}
\end{center}

\begin{enumerate}[leftmargin=1.5em, itemsep=8pt]
  \item Compute $\mu_X$ and $\mu_Y$.

        \noindent\begin{minipage}[t][0.4in][t]{\linewidth}\ans{$\mu_X = 0(0.60)+3(0.30)+10(0.10) = 0.90 + 1.00 = \$1.90$}\\
        \ans{$\mu_Y = 0(0.40)+2(0.40)+6(0.20) = 0.80 + 1.20 = \$2.00$}\end{minipage}

  \item Compute $\sigma_X$ and $\sigma_Y$. (You may use a calculator for the arithmetic.)

        \noindent\begin{minipage}[t][0.5in][t]{\linewidth}\ans{$\sigma^2_X = 0.60(0-1.90)^2+0.30(3-1.90)^2+0.10(10-1.90)^2 = 9.09$; $\sigma_X \approx \$3.02$}\\
        \ans{$\sigma^2_Y = 0.40(0-2)^2+0.40(2-2)^2+0.20(6-2)^2 = 4.80$; $\sigma_Y \approx \$2.19$}\end{minipage}

  \item Which game has the higher expected prize? \ans{Game B (\$2.00 vs.\ \$1.90)}

  \item Which game is riskier (higher $\sigma$)? \ans{Game A (\$3.02 vs.\ \$2.19)}

  \item Define net winnings for each game (prize $-$ \$3). What is the expected
        \emph{net} winnings per play for each game?\\
        \ansline{Game A: $\$1.90 - \$3 = -\$1.10$ per play.}\\[1pt]\ansline{Game B: $\$2.00 - \$3 = -\$1.00$ per play.}

  \item Which game would you choose? Justify using both $\mu$ and $\sigma$.\\
        \ansline{Game B: higher expected net winnings ($-\$1.00$ vs.\ $-\$1.10$) and less}\\[1pt]\ansline{variability ($\sigma_Y \approx \$2.19 < \sigma_X \approx \$3.02$). Both lose money}\\[1pt]\ansline{in the long run, so neither is a good bet.}
\end{enumerate}
\end{tcolorbox}

\end{document}
